% !TeX spellcheck = hu_HU
% !TeX encoding = UTF-8
% !TeX program = xelatex
%----------------------------------------------------------------------------
\chapter{\bevezetes}
%----------------------------------------------------------------------------

Az informatika széles körben alkalmaz különféle nyelveket ahhoz, hogy valamilyen feladatra tudjon megoldást adni:
gyakran ez egy program fejlesztése, ahol egy megfelelően választott  programozási nyelvet (\pl C++, C\#, Java, Rust, Haskell, Erlang, stb.) szokás használni -- a megfelelő választás pedig függ a feladat jellegétől, a projekttel szemben támasztott követelményektől, határidőkről és még magától a fejlesztő csapat előképzettségétől, képességeitől és preferenciájától is.

A problémára nagyszerűen illeszkedő programozási nyelv választásakor is előfordul azonban, hogy fejlesztés közben kerülnek elő olyan részfeladatok, amelyek megoldásához nincs szükség az előbb felsorolt programnyelvek teljes eszközkészletére.
Ilyenkor az is lehetséges, hogy ez az eszköztár akár még meg is nehezítheti a konkrét eset megoldását és rontja annak karbantarthatóságát, mert túl sok lehetőséget nyújt a fejlesztőknek, így egy olyan megoldás áll elő amelyik túlságosan komplex, pedig lett volna lehetőség egy letisztultabb megoldásra is -- ez egyszerűen nem jutott a fejlesztők eszébe a túlságosan tágan hagyott megoldási lehetőségek halmazából.

Ez az egész problémakör azért merül fel, mert ezek a nyelvek \un általános célú programozási nyelvek (\gls{GPL}), amelyek célja és tervezésükkor kitűzött alapvető elvárás, hogy bármilyen feladatra és problémára, amivel szembe találkozhat egy fejlesztő, arra legyen lehetősége azt a nyelven belül megoldani.
Így minden esetben minden probléma megoldásához szükséges minden lehetőséget biztosítanak, akkor is, ha például egy állapotgép leírásához nincs feltétlenül szükségünk arra, hogy tudjunk osztálysablonokat definiálni.

Ez egy gyakran előforduló helyzet, hiszen sok esetben a nagy célprobléma részekre bontása során előállnak olyan részproblémák, amelyek megoldását -- az előző paragrafusban tárgyalt \gls{GPL}-ek túlságosan gazdag lehetőségei miatt -- a nem annyira komplex és kevésbé általános nyelvi keretek könnyítenék.
Ezzel összhangban, sok olyan speciális nyelv alakult ki, amelyik az általános programnyelvek teljes eszköztárát limitált formában, vagy egyáltalán nem biztosítják; cserébe a megcélzott szakterület igényeit elégíti ki a lehető legjobban.
Egy kifejezetten gyakran előkerülő család a leíró nyelvek (\gls{XML} dialektusok, \gls{HTML}, Markdown), de igen nagy számban fordulnak elő különböző egyéb példák: ábrák készítésére készült nyelvek (Pikchr, PlantUML), fordítást kezelő rendszerek (Make, Ninja), de még az adatmanipulációs SQL és QUEL nyelvek és tipográfiai nyelvek, mint a \TeX\ és roff is ebbe a kategóriába tartoznak.

Az előző bekezdésben tárgyalt nyelveket, amelyek egy-egy részprobléma megoldására adnak specializált -- és pusztán arra limitált -- eszközkészletet hívjuk szakterület-specifikus nyelveknek (\gls{DSL}).
Ezek részletesebb tárgyalásáról szól \aref{chr:dsl}.~fejezet.

Az általam fejlesztett általános célú programozási nyelvet -- a C4-et -- és annak általam adott megvalósítását mutatja be \aref{chr:c4}.~fejezet.
Itt ismertetem a tervezési elveket, és azokat az általános elvárásokat, amelyek figyelembe vételével készült a C4 nyelv.
A használt eszközöket felvezetem, és bemutatom a megvalósítás általános struktúráját, kiemelve az érdekesebb, kevésbé elterjedt nyelvi elemekhez tartozó konkrét megvalósítási lépéseket.

\Aref{chr:uses}.~fejezetben a megvalósított nyelvhez mutatok a valóságban is használható felhasználási lehetőséget: megvalósított szakterület-specifikus nyelveket, amelyekben készített kódrészleteket egy az egyben futtatni is lehet a C4 nyelv fordítóját felhasználva.

Ezek után \aref{chr:ftr}.~fejezetben ismertetem a jövőbeli továbbfejlesztési lehetőségeket, és egy olyan nagyobb lélegzetvételű C4 nyelvet felhasználó projektet, amelyik C és C++ forrásfájlok programozható statikus analízisét teszi lehetővé.


















